\chapter{Тёмные пятна}\label{ch:pyatna}

Честный разбор состоит не только из ответов, но и из списка вопросов. Здесь собрано всё, что мы \emph{не} знаем, и причины.

\section{Связь модуля с дроном}

\begin{itemize}
  \item \textbf{Номер метки} (закрыто). С меткой 0 в байтах 2--3 стоял 0, с меткой 13 --- \texttt{00 0d}, то есть 13. Привязка проверена на ненулевой метке. Не проверены только метки больше 255, где старший байт ненулевой (например, 143 помещается в один байт, а 287 нет). \pr{}
  \item \textbf{Другие типы кадров.} Типы 1, 4, 5, \texttt{F9} и кадр \texttt{CE CE} в автопилоте. \nei{}
  \item \textbf{Какие команды дрон шлёт модулю на деле} и с какой частотой. В программе модуля набор команд разобран, живой обмен не снят. \nei{}
  \item \textbf{Как автопилот использует расстояние ToF} (байты 18--19 кадра). \nei{}
  \item \textbf{Чем определяется выбор} между приёмником с одним или двумя байтами заголовка и между вариантами \texttt{int\_dev0} и «ext device uart». Флаг найден, его значение по умолчанию нет. \nei{}
\end{itemize}

\section{Модуль}

\begin{itemize}
  \item Загрузчик и ядро модуля: образ снят целиком, корневая система разобрана (MaixPy3), загрузчик и ядро нет. \nei{}
  \item Размер метки в программе модуля (вывод: около 100 мм) и происхождение масштаба 8,333 в автопилоте. Сама калибровка камеры найдена. \nei{}
  \item Кто на борту штатно использует порт UART0 и команду \texttt{FA}. Сам канал работает (зуммер). \nei{}
  \item Формат \texttt{color\_config.json} (на модуле просто \texttt{null}); есть ли подпись у обновления (видна лишь контрольная сумма). \nei{}
  \item Совпадает ли эта программа с камерой C10. \nei{}
\end{itemize}

\section{Автопилот}

\begin{itemize}
  \item Константы регуляторов, драйверы датчиков и исполнитель миссии известны по названиям и по ядру C10, но не по самому \texttt{pilot\_f}. \nei{}
  \item Закрытый алгоритм оптического потока (\texttt{flow}). \nei{}
  \item Подпись образа дрона, возможность вернуть прежнюю прошивку. \nei{}
  \item Какая форма команд принимается вашим Wing по UDP 9696: подгруппа~3 или подгруппы 12--14. По USB-пульту команды скорости работают (глава~\ref{ch:pult}). \nei{}
  \item Кто и по какому событию переписывает \texttt{/netPrivate/wifiActor}; зачем нужен TCP-порт 9697. \nei{}
  \item Остальные команды зуммера (\texttt{02}--\texttt{0C}), громкость и остановка ноты. \nei{}
  \item Как запускать ретранслятор для зуммера при включении дрона: нужна запись во флеш (\texttt{autorun.sh}), мы её не делали. \nei{}
  \item Положение, посчитанное дроном: проверено вживую на метке 13, формула верна (глава~6). Не проверено на других высотах и при наклоне. \pr{}
\end{itemize}

\section{Протокол приложений}

\begin{itemize}
  \item Единицы телеметрии (напряжение, ёмкость, углы). \nei{}
  \item «Версия» в байте 57 кадра по трассировке соответствует другому полю. \nei{}
  \item Порядок слотов моторов в команде \texttt{getEngine} (мотор 1, вероятно, в последнем слоте) и смысл \texttt{getAexsSelect}. \nei{}
\end{itemize}

\section{Совсем не разобрано}

Прошивка пульта (GD32E103), Linux-ядро C10 (V853S) и его камера, схемы плат, страницы документации производителя (не загрузились).

\section{Что дальше даст больше всего}

\begin{enumerate}
  \item Снять живой обмен: что и как часто дрон шлёт модулю.
  \item Проверить, переживает ли ретранслятор отключение кабеля ADB; тогда шоу с поющим дроном получится без записи во флеш.
  \item Запустить свой \texttt{main.py} вместо штатной программы модуля (на столе, без пропеллеров) и отправить автопилоту кадр с другой меткой.
  \item Проверить положение на нескольких высотах и при наклоне, а не на одной точке.
  \item Разобрать приёмник автопилота \texttt{pilot} в режиме точки доступа: он отличается от \texttt{pilot\_f}.
\end{enumerate}
